% \chapter[测量]{\itr{Measurement}{测量}}
\chapter[基础知识]{\itr{Basic Knowledge}{基础知识}}
本章节内容较为简单，一般相关题目会出现在日常上课的课后小测与作业内，大型考试基本不涉及。重点是掌握物理学的基本单位，以及利用国际单位制以及相关比例关系推导特殊物理量的单位表达式与可能的公式表达式。
\section[单位与量纲]{\itr{Unit \& Dimension}{单位与量纲}}
\begin{Itemize}
    \item \itr{Basic Quantity}{基本量}\ \eg Length($L$), Mass($M$), Time($T$)\\
    在自然世界中，我们会用各种物理量来描述物质的性质。就像在平面几何中我们可以基于几大公理推出各种定理，在物理中，我们也可以规定一些“基本的量”，使得所有的物理量可以由这些量导出。自然地，我们就称呼这些量为基本量。
    \item \itr{Unit}{单位}在规定好一些物理量之后，我们会想办法去度量它们。既出于定量测量的要求，又出于减少物理公式参数的考量，人们定义了各种各样的单位。
    \begin{itemize}
        \item \itr{SI Units}{国际单位制单位}\mgnote{SI $\Leftrightarrow$ International System of Units}\ \eg meter(m), kilogram(kg), second(s)
        \item \itr{Non SI Units}{非国际单位制单位}\ \eg mm, $\upmu$m, g\\
              \En{If you \itr{convert}{转化} everything to the basic SI units, you can \itr{omit}{省略} them during the calculation, but remember to put back the correct units at the end.}\\
              事实上，可以在计算过程中省略的约定也算是国际单位制存在的意义之一(\dove :大概吧)。
    \end{itemize}
    \item \itr{Dimension}{量纲}\ 一个物理量的量纲指的是这个物理量关于基本量的导出式。

    \eg 速度的量纲可以表示为[$v$]=$\dfrac{L}{T}$

    \item \itr{Dimension Analysis}{量纲分析}\ 量纲分析是检验物理公式正确性的好方法。简而言之，你可以简单地通过分析一个物理等式左右两边的量纲是否一致来初步判断这个等式可不可能成立。当然，量纲分析也可以作为导出实验公式大方向的一个指导，又或者……是你忘记了公式中的某一个物理量时尝试硬添物理量时的救命稻草。
\end{Itemize}
\section[数据处理]{\itr{Data Processing}{数据处理}}
这部分普通物理学实验\Romannumeral{1}\footnote{说到这里，\dove 还编写了一个用于普物实验数据处理的软件，链接为 \url{https://github.com/CrazySpottedDove/Lab-Assistance.git}，如有需要，可以自行下载。}会详细说明并实际计算运用，非常抱歉，此处略去。

\section[微分方程]{微分方程}
在普通物理学中，一个非常重要的思想便是微分的思想，也间接导致我们我们有时无法直接写出具体的公式解析式，而是必须先通过微元法得到微分方程，再行进行数学求解。以往的授课中对于微分方程列写以及求解的方法大多通过背公式的方式求解微分方程。但事实上，部分公式完全可以通过现推的方式解决，省去一部分的记忆负担。

本节我们将基于工科常微分方程的内容进行讲解，对于微分方程解的存在性这类分析学的问题，或者更难的偏微分方程，由于课程并不涉及，因此本节将略去。若平时或考试中遇到偏微分方程等目前无法解决的微分方程，请先确认是否存在满足可分离变量得到常微分方程的形式，或是否方程列写错误。

\subsection[微分方程引入]{\itr{Introduction to Differential Equations}{微分方程引入}}
微分方程，顾名思义，就是包含自变量及其导数组成的方程。例如对于自变量$x$，因变量$y=f(x)$，$y' = x$就是一个经典的微分方程。进一步地，微分方程可以有如下几种简单的分类：
\begin{Itemize}
    \item 根据自变量的个数，分为常微分方程和偏微分方程。例如，如果因变量$y$是只与$x$有关的函数，例如我们前面的例子，那就是一个常微分方程；如果$y$是一个多元函数（间接地其导数需要写作偏导数的形式），那就是一个偏微分方程。
    \item 根据方程中自变量及其导数的系数，分为线性方程与非线性方程。如果方程满足线性性，那么就是线性方程，反之为非线性方程。一般的，系数如果是与自变量无关的常数，那么就是线性的。
\end{Itemize}

就该课程来说，我们只需要知道是或者不是微分方程，重点在于如何求解。微分方程的求解结果不是具体的自变量的值，而是未知函数的函数表达式。

为了求解微分方程，我们首先需要列写微分方程。列写过程与物理中的微元法有异曲同工之妙。在中学的微元法中，我们通常假设一个很小的时间间隔$\Delta t$，或者一个很小的位移$\Delta x$，然后把他们写作一个求和的形式。但是从积分的定义理解，无穷个无穷下的求和应该可以写作积分的形式，积分的过程也是沿着$x$轴方向切分成无数的矩形后求和，即：
\[\int_a^b f(x)\dif{x} = \lim_{\Delta x \to 0} \sum f(x_i)\Delta x\]

这时不难理解，$\dif{x}$同时也表示了一种微元的概念，而积分符号则是一种无穷求和的概念，$f(x)\dif{x}$表示的就是这个小矩形的面积。因此这样的微元也可以被赋予一种物理意义。例如很小的一段长度的微元是$\dif{x}$，很小一块体积的微元的质量为$\dif m = \rho \dif{V}$，很短时间内物体的位移的微元为$\dif{x} = v\dif{t}$。这里的微分符号部分可以表示具体物理意义的变量，并应用物理公式表示该微元物体的相关物理量。在列写出相关方程后，微分符号则不加变动地具备指明被积参数的数学含义。

\subsection[一阶微分方程求解]{一阶微分方程求解}
一阶微分方程，即最高阶导数仅包含到一阶导数的微分方程，其通式可以写作：
\[f(x,y') = 0\]

\begin{enumerate}
    \item 一阶可分离变量方程，形如：
    \[\dfrac{\dif y}{\dif x} = f(x)g(y)\]

    求解该方程，只需要将$x$和$y$完全分离到等号两边，并做积分即可。积分的形式如下：
    \[\int \dfrac{\dif y}{g(y)} = \int f(x)\dif x\]

    一般的不定函数积分，我们需要小心积分结果的$+C$。对于实际数值的积分，积分上下限要求两边是对应的。如果$x:a\rightarrow b$，那么$y:y(a)\rightarrow y(b)$。另外，对于$g(y) = 0$的情况，我们需要进一步进行讨论，以判断是否有诸如$y = C$的形式的解存在。

    \item 一阶齐次方程
    \[\dfrac{\dif y}{\dif x} = f(\dfrac{y}{x})\]

    求解该方程，需要令$u = \dfrac{y}{x}\quad u = \phi(x)$，也就是$y = ux$。代回微分方程，根据函数乘积的微分可以得到：
    \[\dfrac{\dif y}{\dif x} = u + x\dfrac{\dif u}{\dif x} = f(u)\]

    这可以化为前面的一阶可分离变量方程，即：
    \[\dfrac{\dif u}{f(u) - u} = \dfrac{\dif x}{x}\]

    根据之前的过程计算即可。同样地，要注意分母为0的情况，否则分离变量过程涉及除法计算不合法。

    \item 一阶线性方程
    \[\dfrac{\dif y}{\dif x} + p(x)y = f(x)\]

    根据右边$f(x)$，若其为0，则为一阶齐次方程，否则为非齐次方程。无论哪种类型，都从齐次方程入手。根据可分离变量的原则，可以首先求到：
    \[y = C_1 e^{-\int p(x)\dif x}\]

    在非齐次情况下，这里的常数$C_1$也是自变量的函数，我们可以写作$u(x)$，因此：
    \[y = u(x)e^{-\int p(x)\dif x}\]

    代回之前的微分方程，有：
    \[u'e^{-\int p(x)\dif x} + u e^{-\int p(x)\dif x}(-p(x)) + p(x) u e^{-\int p(x)\dif x} = f(x)\]

    可解得：
    \[u = \int f(x) e^{\int p(x)\dif x}\dif x\]

    因此其完整的解为：
    \[y = \left[\int f(x) e^{\int p(x)\dif x}\dif x + C\right] e^{-\int p(x)\dif x}\]

    \item 一阶线性伯努利方程
    \[\dfrac{\dif{y}}{\dif{x}} + p(x)y = f(x)y^{n}\quad n\neq 1\]

    两边同除$y^n$，令$z=y^{1-n}$，则原方程可化为
    \[\dfrac{1}{1-n}\dfrac{\dif{z}}{\dif{x}} + p(x)z = f(x)\]
    
    使用上面提到的方法继续求解即可。
    \item 可降解为一阶的二阶方程
    \begin{itemize}
        \item 没有一阶微分项：
        \[\dfrac{\dif^2{y}}{\dif{x^2}}=f(x)\]
        连续进行两次积分即可
        \item 没有$y$的常数项：
        \[\dfrac{\dif^2{y}}{\dif{x^2}}=f(x,\dfrac{\dif{y}}{\dif{x}})\]
        令$u=\dfrac{\dif{y}}{\dif{x}}$，则方程可顺利降为一阶方程$\dfrac{\dif{u}}{\dif{x}}=f(x,u)$
        \item 没有$x$的常数项：
        \[\dfrac{\dif{y}}{\dif{x}}=f(y,\dfrac{\dif{y}}{\dif{x}})\]
        令$u=\dfrac{\dif{y}}{\dif{x}}$，则原方程可化为$u\dfrac{\dif{u}}{\dif{y}}=f(y,u)$
    \end{itemize}
\end{enumerate}

\subsection[高阶常系数微分方程求解]{高阶常系数微分方程求解}
线性的高阶微分方程看似更难，实则更简单。对于方程
\[a_0 y + a_1 y' + a_2 y'' + \cdots + a_n y^{(n)} = f(x)\]

其分为齐次解和特解两部分。齐次解是方程右边为0时的解，可以理解为是否包含这些项不影响最终的微分方程的表达式，好比是一个正常人都有五官。特解是满足$f(x)$的解。一般而言，特解部分与$f(x)$具备类似的函数表达式，例如$f(x)$是多项式时特解具备多项式的形式，$f(x)$是正弦函数时特解也包含同频率的正弦函数的形式。特解的作用在于找到特殊的函数以符合方程右边，好比是每个人特殊的五官构型，以区分开不同的人以免认错\footnote{在控制学科中，系统结构可以用微分方程表述。对于一个输入，其齐次解部分告诉我们输出响应瞬态如何，是单调变化还是往复振动；特解部分告诉我们系统稳定之后的输出响应如何，是正弦波动还是常值稳定亦或者多项式形式变化等}。

对于特解，我们一般使用待定系数法，假设一个类似的函数表达式，代入微分方程求解待定系数即可。因此我们重点找齐次解。这里我们使用特征方程，对于上面的微分方程，其特征方程形式如下：
\[a_0 + a_1\lambda + a_2\lambda^2 + \cdots + a_n\lambda^n = 0\]

由此我们可解得这个多项式方程中的特征根$\lambda$的值以及其重数。
\begin{Itemize}
    \item 对于任意一个n重实根$\lambda_i$，其对应齐次解为$e^{\lambda_i x}g_i(x)$，其中$g_i(x)$为最高项$n-1$次的多项式，多项式系数由待定系数确定。
    \item 对于任意一组n重共轭复根$a_i\pm b_i i$则对应表达式$e^{a_i x}[h_i(x)\cos(b_i x)+w_i(x)\sin(b_i x)]$，其中$h_i(x)$和$w_i(x)$都是最高项为$n-1$次的多项式，多项式系数由待定系数确定。
\end{Itemize}

对于含有多个特征根的的常系数线性微分方程，其最终结果为所有特征根对应的齐次解的表达式以及特解表达式之和，其中含有的所有系数为待定系数。若给定额外的初值条件，则可以通过待定系数求解。

我们这里举个例子，对于微分方程：
\[y''+2y'+y=0\]

其特征根为2重实根-1，那么齐次解即为：$y=e^{-x}(c_1+c_2x)$，代回原方程最终求解待定系数即可。

更为一般地，我们经常见到的是二阶微分方程，例如普物中的振动方程就是一个经典的二阶方程。该方程的一般解法与高阶微分方程的求解一致，详细分为下列三种情况，并具备对应的物理意义：
\begin{itemize}
    \item 特征方程为两个不同实根：
    假设特征根为$\lambda_1,\lambda_2$，则齐次解为：
    \[y=c_1e^{\lambda_1x}+c_2e^{\lambda_2x}\]
    \item 特征方程有两个相同实根：
    假设特征根为$\lambda_1=\lambda_2=\lambda$，则齐次解为：
    \[y=(c_1+c_2x)e^{\lambda x}\]
    \item 特征方程为一组共轭复根：
    假设特征根为$\lambda_1,\lambda_2=a\pm bi$，则齐次解为：
    \[y=e^{ax}[c_1\cos(bx)+c_2\sin(bx)]\]

    一般三角函数会合并，并辅以相位参数$\phi$。
\end{itemize}

在物理上例如振动、衰变过程的特征根（或者特征根的实部）是负数，因此齐次解是指数衰减的方程，上面情况1对应方程没有三角函数项，呈现单调性，称为过阻尼；情况3具备反复震荡的三角函数部分，称为欠阻尼；情况2的特征根介于1、3之间的临界状态，称为临界阻尼。

欠阻尼的一个实际例子为受到空气阻力的单摆的摆动方程，当然这个例子没有那么严谨，但我们可以大致感受到其实际现象。
\section[三角函数、复数与其矢量表示]{三角函数、复数与其矢量表示}
提到这一部分知识，最经典的应该是从欧拉公式起步。我们知道很经典的等式：
\[e^{i\pi} + 1 = 0\]

事实上，欧拉公式表示为：
\[e^{i\theta} = \cos\theta + i\sin\theta\]

一种偏计算法的证明方式为：对于指数函数的泰勒展开，代入自变量$i\theta$，并对该级数进行重排，反向利用三角函数的泰勒展开得到最终结论：
\begin{align*}
    e^{i\theta} &= 1 + i\theta - \dfrac{\theta^2}{2!} - \dfrac{i\theta^3}{3!} + \dfrac{\theta^4}{4!} + \cdots\\
    &= (1 - \dfrac{\theta^2}{2!} + \dfrac{\theta^4}{4!} + \cdots)\\
    &+ i(\theta - \dfrac{\theta^3}{3!} + \dfrac{\theta^5}{5!} + \cdots)\\
    &= \cos\theta + i\sin\theta
\end{align*}

那么现在我们知道，如果要取余弦函数，就取该复指数的实部，正弦函数则取虚部。因此，三角函数可以采用复指数的方式计算，最终取其中某个方向的分量即可。

我们再观察一下上面的等式的形式。两个三角函数的存在让我们联想到极坐标，同时考虑到复数在数轴上的表示采用了复平面，那么我们可以很自然地联系在一起。以实轴为$x$轴，虚轴为$y$轴，那么一个复数可以在复平面上找到唯一的一个点。

根据欧拉公式的结论，我们可以构建这个点与原点（即表示数字0）的点之间的连线，这个时候自然会直观地显示出我们的角度$\theta$。下图表示了这一过程，习惯上，我们的角度从$x$轴即实轴正半轴出发，逆时针方向旋转为正角度。\footnote{这种角度正负号的定义广泛应用于各种与平面坐标系相关的角度约定}

\begin{singlefigure}[复数在复平面上的表示]{chapter1_复数在复平面上的表示}[0.5]
\end{singlefigure}

而在二维平面上的有向箭头，又可以表示为一种矢量。复数可以表示实数加虚数，也可以表示成复指数，前者适合加减法运算，而后者适合乘除法运算，给我们的计算带来了很大的便利。这样我们就把复数与矢量构建了一个初步的认知，而三角函数与又可以用复数表示，这或许可以带给我们一些启发。

上面的结论主要用于物理中的各种分析和计算过程，在现实世界中，三角函数是直观展现在我们面前的，但复指数和矢量提供了更高维度的视角去观察三角函数。当我们考虑一个波动的时候，拉展开时间线，一个一维运动的波动就如一个沿着时间轴展开的螺旋线在平面上的投影。物理世界中，各种波动和振动都是与时间有关的两，频率与初相位是其重要参数。假定两个典型的波动方程：
\[y_1 = A_1 \cos (\omega_1 t + \phi_1)\qquad y_2 = A_2 \cos (\omega_2 t + \phi_2)\]

我们有时会考虑两个波的超前滞后性。当两个波动的频率相等时，两个波动的相对相位关系将会被固定，如果我们略去共同的频率和时间项，仅保留幅值和初相位，一个是大小一个是角度，自然我们会发现和极坐标以及复平面上的复数非常相似，我们可以得到下面的示例：
\begin{singlefigure}[相量表示的波的相对关系]{chapter1_相量表示的波的相对关系}[0.7]
\end{singlefigure}

以幅值为长度，初相位为从$x$正半轴逆时针转过的角度，可以用相量（可以理解为一种矢量，一些地方也称为相幅矢量）表示出两个波。在整个时间变化过程中，这部分整体会不断地旋转，但相对位置关系保持不变。这里的矢量也被称为旋转平面矢量，它是一种时间矢量，因为我们要求波动的频率相同，因此它包含了描述一个波动所需的大小和初始相位两个其余指标。我们一般取$x$轴正方向为参考相量。旋转矢量的大小既可以用振幅表示，也可以用有效值表示，只要同一个相量图中采用同样类型的值，计算便是有效的。

既然是一个矢量，那么它的加减法便遵循平行四边形法则。上图所示中我们绘制了两个矢量的和矢量，得到的矢量可以得知它的相对于参考相量的相位以及它的大小，那么我们便可以得到两个波动函数的加和对应的波动方程的大小和初相位。当然我们可以再举一个例子：
\[y_1 = \cos(\omega t)\qquad y_2 = \sqrt{3}\cos(\omega t+90^\circ)\]

如果要计算三角函数的求和，单纯相加会相对复杂，但是如果采用相量的方法，我们绘制出矢量表示的波，采用平行四边形加法可以很明显地得到：
\[y = y_1 + y_2 = 2\cos(\omega t + 60^\circ)\]

当然，相量的作用不仅仅是面向波的加减法，也能够便于微积分的计算。当我们把三角函数部分统一放在复指数的指数部分时，对该三角函数的求导将会等价于对复指数的求导，例如：
\[y = A_1\cos(\omega t + \phi)\quad\Leftrightarrow\quad e^{i(\omega t + \phi)}\]
\[y' = A_1\cos(\omega t + \phi + 90^\circ)\quad\Leftrightarrow\quad i\omega e^{i(\omega t + \phi)}\]

而在复平面上，一个复数乘$i$等价于这个复数对应矢量逆时针旋转$90^\circ$，因此结果上依然是等价的，复指数的求导结果对应的旋转矢量依然等价于三角函数的结果，该三角函数被求导后依然是这个被求导后的复指数的实部。

在我们前面讨论的过程中我们可以了解到复指数表示下的三角函数的表示以及这种表示对于计算等过程的极大简化。但我们必须提醒，单一复指数表示的三角函数与原三角函数并不等同。在一般的线性运算中（如加减法或求导和积分），复指数将三角函数的计算转化为矢量的旋转和加和，其中加减法过程要求不同三角函数中的频率$\Omega$一致才可以用相量计算。而在乘除法等过程中，如果单一地将三角函数括号内的部分转移到指数部分中构建复指数，将会推导出错误的结果。事实上，根据欧拉方程，三角函数的标准真实值应当表示为：
\[\cos\theta = \dfrac{e^{i\theta} + e^{-i\theta}}{2}\qquad \sin\theta = \dfrac{e^{i\theta} - e^{-i\theta}}{2i}\]

因此在使用时首先要考虑是否满足前提条件。一般而言，复指数和相量法适用于振动与波（以及采用振动与波表示的一些物理场景，例如波动光学）以及含容含感电路（电容电感引入了求导和积分过程）的电路计算，将三角函数转化为矢量计算、将微积分过程转化为四则运算的方法让我们得以用几何直观解释三角函数的过程，也从更高的维度观察到振动的表示。